\documentclass[11pt]{article}

%% ---------------------------------------------------------------
%% Packages
%% ---------------------------------------------------------------
\usepackage[margin=1in]{geometry}
\usepackage{graphicx}
\usepackage{booktabs}
\usepackage{array}
\usepackage{multirow}
\usepackage{amsmath}
\usepackage{amssymb}
\usepackage{bm}
\usepackage{authblk}
\usepackage[hidelinks]{hyperref}
\usepackage{xcolor}
\usepackage{caption}
\usepackage{subcaption}
\usepackage{url}
\usepackage{microtype}
\usepackage{enumitem}
\usepackage[numbers,sort&compress]{natbib}
\usepackage{parskip}
\usepackage{setspace}
\usepackage{mdframed}
\usepackage{ifthen}

%% ---------------------------------------------------------------
%% Hyperref setup
%% ---------------------------------------------------------------
\hypersetup{
    colorlinks=true,
    linkcolor=blue!70!black,
    citecolor=blue!70!black,
    urlcolor=blue!70!black
}

%% ---------------------------------------------------------------
%% Figure placeholder command (for schematics not yet drawn)
%% ---------------------------------------------------------------
\newmdenv[
  linecolor=gray!60,
  backgroundcolor=gray!8,
  linewidth=1pt,
  innertopmargin=10pt,
  innerbottommargin=10pt,
  innerrightmargin=10pt,
  innerleftmargin=10pt,
  skipabove=6pt,
  skipbelow=6pt
]{figframe}

%% Usage: \figplaceholder{short title}{height}{description}
\newcommand{\figplaceholder}[3]{%
  \begin{figframe}
    \centering
    \rule{\linewidth}{0pt}\\[#2]
    {\small\textbf{[Figure: #1]}\\[4pt]
    \textit{#3}}
    \rule{\linewidth}{0pt}
  \end{figframe}%
}

%% Safe include: show figure if it exists, otherwise show a placeholder.
%% Note: \IfFileExists does not honor \graphicspath, so we probe the
%% figures/ subfolder first and fall back to the current directory.
\newcommand{\safeinclude}[2][width=\textwidth]{%
  \IfFileExists{figures/#2}{%
    \includegraphics[#1]{figures/#2}%
  }{%
    \IfFileExists{#2}{%
      \includegraphics[#1]{#2}%
    }{%
      \figplaceholder{File not yet generated}{2.5cm}{Path: \texttt{\detokenize{#2}}. Generate this figure and recompile.}%
    }%
  }%
}

%% ---------------------------------------------------------------
%% Typography
%% ---------------------------------------------------------------
\setstretch{1.15}
\setlength{\parindent}{0pt}
\setlength{\parskip}{6pt}
\captionsetup{font=small, labelfont=bf, width=0.95\linewidth}

%% Figure directory
\graphicspath{{figures/}{./}}

%% ---------------------------------------------------------------
%% Title and authors
%% ---------------------------------------------------------------
\title{\textbf{Radiomics-Based Machine Learning and Deep Learning
for Clinically Significant Prostate Cancer Detection
on Biparametric MRI}}

\author[1,*]{Jesus Alejandro Alzate-Grisales}
\author[4,6]{Mario Alejandro Bravo-Ortiz}
\author[4,9]{Luis Miguel Cardona Pérez }
\author[4]{Salomé López Marulanda}
\author[2]{Miguel Perán-Teruel}
\author[2]{Azahar Navarro Beltrán}
\author[2]{Clara Ruiz Torres}
\author[2]{José Manuel Osca García}
\author[3]{Francisco García-García}
\author[4,5,6,7]{Reinel Tabares-Soto}
\author[8]{Juan M García-Gómez}
\author[1,*]{Maria de la Iglesia-Vayá}
\affil[1]{Unidad Mixta de Imagen Biomédica e Inteligencia Artificial FISABIO-CIPF, Fundación para el Fomento de la Investigación Sanitario y Biomédica de la Comunidad Valenciana, Valencia 46020, Spain}
\affil[2]{Hospital Arnau de Vilanova, Urology, Valencia, 46015, Spain}
\affil[3]{Computational Biomedicine Laboratory, Principe Felipe Reseach Center (CIPF), 46012, Valencia, Spain}
\affil[4]{Electronics and Automation Department, Universidad Autónoma de Manizales, Manizales, 170001, Caldas, Colombia}
\affil[5]{Faculty of Engineering and Sciences, Universidad Adolfo Ibáñez, 7941169, Santiago, Chile}
\affil[6]{Department of Systems and Informatics, Universidad de Caldas, Manizales, 170001, Caldas, Colombia}
\affil[7]{GobLab School of Government, Universidad Adolfo Ibáñez, Santiago, Chile}
\affil[8]{BDSLab, Instituto Universitario de Tecnologías de la Información y Comunicaciones, Universitat Politècnica de Valencia, 46022 Valencia, Spain.}
\affil[9]{Centro de Bioinformática y Biología Computacional de Colombia, BIOS. Manizales, Colombia}
\affil[*]{jesus.alzate@fisabio.es, delaiglesia\_mar@gva.es}

\date{}

%% ===============================================================
\begin{document}
\maketitle

%% ---------------------------------------------------------------
\begin{abstract}
\noindent
We benchmarked radiomics-based classical machine learning, deep tabular
learning, and a pretrained tabular foundation model (TabFM) for case-level
detection of clinically significant prostate cancer (csPCa) on biparametric
MRI, and assessed the added value of routine clinical variables. The full
PI-CAI Public Training and Development Dataset (1{,}500 studies, 1{,}476
patients, 425 csPCa-positive, 28.3\% prevalence) was analyzed under the
official patient-grouped five-fold cross-validation. Whole-gland radiomic
features from T2-weighted, apparent diffusion coefficient, and high b-value
diffusion-weighted imaging were used to train three classical models (Random
Forest, Gradient Boosting, LightGBM), three deep models (tabular Transformer,
CapsNet, Transformer--CapsNet hybrid), and TabFM on a shared, leakage-safe,
fold-specific feature-selection plan with calibration and an
outer-training-derived threshold. Age, prostate-specific antigen (PSA), PSA
density, and prostate volume were evaluated alone, by concatenation, and by
dual-branch (or, for TabFM, dual-stream late-fusion) combination. Metrics were
computed on aligned out-of-fold predictions with patient-level cluster-bootstrap
95\% confidence intervals. TabFM gave the best radiomics-only discrimination
(AUROC 0.828, 95\% CI 0.805--0.850; AUPRC 0.650), numerically above the Random
Forest (AUROC 0.814) and the deep tabular pipelines (AUROC 0.794--0.803),
although the TabFM-versus-Random-Forest difference was not significant after
multiplicity adjustment ($\Delta$AUROC 0.014; Holm-adjusted $P=0.084$).
Radiomics clearly outperformed clinical variables alone (best clinical-only
AUROC 0.757, also TabFM); combining the two sources produced the best models
overall, whether by simple concatenation passed to TabFM (AUROC 0.844, 95\% CI
0.822--0.865; AUPRC 0.666) or by a logistic late-fusion of separately trained
radiomics-only and clinical-only TabFM streams (AUROC 0.844, 0.822--0.865;
AUPRC 0.674); both were statistically indistinguishable from each other
($P=1.00$) but each was a significant improvement over the previous best model
in its respective condition (concatenated Random Forest, Holm-adjusted
$P=0.020$; dual-branch Transformer, Holm-adjusted $P=0.004$), and the
concatenated TabFM was also a significant improvement over the radiomics-only
Random Forest (Holm-adjusted $P<0.001$). A
random-weight control, in which TabFM's pretrained parameters were replaced by
random initialization while keeping the architecture and inputs unchanged,
collapsed to near-chance discrimination (AUROC 0.522), indicating that the
observed gain reflects the foundation model's pretraining rather than its
architecture alone. Center-stratified analysis revealed marked heterogeneity,
with significantly lower discrimination at one of three centers (AUROC 0.74
vs.\ 0.84 elsewhere) in every model family including TabFM. Under leave-one-center-out
retraining---a direct test of transportability to an unseen center---discrimination fell
substantially for every model family (radiomics-only pooled AUROC $\approx$0.68--0.70,
combined $\approx$0.71--0.75), with TabFM retaining a small numerical lead
(radiomics-only 0.702; combined 0.753) but no model reaching its internal
cross-validation performance, quantifying an optimistic bias of $\approx$0.10--0.14 AUROC
in the center-mixed folds. Model-agnostic
SHAP attributions confirmed that TabFM relied on the same diffusion-weighted
and PSA-density signature as the classical and deep models. A pretrained tabular
foundation model thus led numerically in all conditions and provided a
modest but statistically supported gain over strong, calibrated classical and
deep tabular baselines in the combined-modality conditions, with a consistent
but non-significant lead in the radiomics-only and clinical-only settings, while
a calibrated classical ensemble remained a highly competitive and more
parsimonious alternative; aggregate
performance again masked clinically relevant between-center variation. These
comparisons are exploratory, as model families were screened on the same cohort.
\end{abstract}

%% ---------------------------------------------------------------


%% ---------------------------------------------------------------
\section{Introduction}
\label{sec:intro}

Prostate cancer (PCa) is one of the most frequently diagnosed malignancies in men and remains a leading cause of cancer-related mortality, making accurate detection and risk stratification essential for guiding management~\citep{Midiri2021Diagnostics,li2022machine}. Multiparametric magnetic resonance imaging (mpMRI) has become a key tool for the non-invasive evaluation of suspected PCa and is now recommended before biopsy in many diagnostic pathways. To this end, the Prostate Imaging Reporting and Data System (PI-RADS) provides a standardized framework for image acquisition, interpretation, and reporting~\citep{Midiri2021Diagnostics,li2022machine}. However, PI-RADS assessments remain subject to inter-reader variability and may lead both to missed clinically significant disease and to unnecessary biopsies in men with equivocal lesions~\citep{Midiri2021Diagnostics}.

Radiomics has emerged as a promising approach to address these limitations by extracting large numbers of quantitative features from medical images and using them as objective imaging biomarkers for tumor detection, grading, and prognosis~\citep{Midiri2021Diagnostics,Bao2024}. Recent narrative and systematic reviews of MRI-derived radiomics models for PCa highlight applications in lesion detection, Gleason score prediction, and outcome estimation. These reviews report that well-constructed models often achieve areas under the receiver operating characteristic curve (AUROC) in the range of approximately 0.80 to 0.90 during internal validation; however, they also underscore substantial methodological heterogeneity, emphasizing the need for standardized pipelines and transparent reporting~\citep{Midiri2021Diagnostics,JZUSciB2023MRIRadiomics,collins2024tripodai}.

Radiomics-based machine-learning (ML) models have been developed on both mpMRI and biparametric MRI (bpMRI), typically utilizing T2-weighted and diffusion-weighted or apparent diffusion coefficient (ADC) sequences, with lesion- or gland-level segmentations serving as regions of interest~\citep{QIMS2022PerformanceVar,Diagnostics2022BpMRIReview}. In heterogeneous multicenter datasets, tree-based classifiers—such as Random Forests—trained on engineered shape, intensity, and texture features have achieved AUROCs of 0.78 to 0.83 for clinically significant PCa (csPCa). When appropriate bias-mitigation and cross-validation strategies are applied, these models perform comparably to conventional imaging biomarkers such as mean ADC, PI-RADS, and prostate-specific antigen (PSA) density~\citep{QIMS2022PerformanceVar,Diagnostics2022BpMRIReview}. More recently, a multicenter retrospective study in real-world clinical practice evaluating 1{,}616 patients across four tertiary centers demonstrated that Random Forest-based models reached AUROCs of 0.874 to 0.893 across internal and external test cohorts, and that radiomics-adjusted PI-RADS scores improved specificity with only a slight sacrifice to sensitivity~\citep{Bao2024}.

The relative performance of classical radiomics-based ML versus deep learning (DL) remains an active topic of investigation. A validation study by Castillo et al.\ compared a radiomics model based on delineated mpMRI regions with a fully automated DL model for csPCa classification across three independent test cohorts. While DL obtained a higher AUROC during internal cross-validation, the classical radiomics model achieved a higher AUROC on all external cohorts, suggesting better generalizability of the radiomics approach under those conditions~\citep{Castillo2022Cancers}. Systematic reviews of bpMRI further support that while both paradigms can reach AUCs in the mid-0.80 range, DL often requires substantially larger training cohorts and may be particularly sensitive to domain shift, whereas tree-based ensembles remain highly competitive on tabular radiomic feature spaces~\citep{Diagnostics2022BpMRIReview,castillo2020automated}.

A distinct and rapidly emerging paradigm for tabular data is the tabular foundation model, which reframes prediction as in-context learning: rather than fitting weights to each new dataset, a single Transformer pretrained on large collections of synthetic tabular tasks infers the mapping between features and labels directly from the training rows provided at inference time. TabPFN demonstrated that this approach can outperform tuned tree ensembles and neural networks on small-to-medium tables while requiring no dataset-specific training~\citep{hollmann2025tabpfn}, and subsequent models have extended in-context learning to larger tables~\citep{qu2025tabicl}. The recently released TabFM further scales this strategy through pretraining on hundreds of millions of synthetic datasets~\citep{kong2026tabfm}. Because engineered radiomic feature tables are moderate-dimensional, strongly inter-correlated, and typically available only in the low-thousands of cases, they fall squarely in the regime where such models are expected to be competitive, yet tabular foundation models have not, to our knowledge, been benchmarked for csPCa detection against calibrated classical and deep tabular baselines on an identical, leakage-safe radiomic feature space.

An important line of research has focused on combining these radiomics profiles with routine clinical variables to build multimodal risk models. In a large single-center retrospective study (1{,}497 MRI exams from 1{,}395 men), Antol{\'i}n et al.\ developed four Random Forest-based models using automatically segmented whole-gland radiomics, PI-RADS, and clinical variables~\citep{Antolin2026Insights}. In their validation cohort, radiomics alone did not significantly outperform PI-RADS (AUC 0.838 vs.\ 0.833), but the combined clinical-radiomics-PI-RADS model achieved an AUC of 0.891, significantly outperforming the simpler models and providing the optimal trade-off between sensitivity, specificity, and biopsy avoidance~\citep{Antolin2026Insights}. Other clinical-radiomics models targeting PI-RADS~3 lesions or transition-zone tumors likewise confirm that adding radiomics to routine predictors—such as age, PSA, PSA density, and prostate volume—can yield modest but clinically meaningful gains in discrimination and decision-curve net benefit~\citep{BMC2025PIRADS3,FrontOnc2026TZRadiomics}.

Against this background, large multicenter datasets and standardized evaluation protocols have been introduced to benchmark AI models for csPCa detection. Recent multicenter studies have evaluated MRI-based models designed to mimic routine practice, demonstrating that calibrated radiomics models can lower the proportion of equivocal cases while maintaining high negative predictive values~\citep{Bao2024,JZUSciB2023MRIRadiomics}. Concurrently, methodological reviews and the TRIPOD+AI statement call for more rigorous evaluation protocols, including patient-grouped cross-validation, strictly isolated feature-selection strategies, probability calibration, decision-curve analysis, and formal statistical comparison of competing models~\citep{collins2024tripodai}.

The present study addresses several of these methodological gaps by benchmarking three families of tabular predictors---classical machine learning, deep tabular learning trained from scratch, and a pretrained tabular foundation model (TabFM)---on an identical, IBSI-compliant whole-gland radiomic feature space derived from bpMRI. This evaluation utilizes the large, multicenter Prostate Imaging: Cancer AI (PI-CAI) public training and development dataset (1{,}500 studies from 1{,}476 patients, csPCa prevalence 28.3\,\%). Radiomic features are extracted from T2-weighted imaging, ADC maps, and high-$b$-value diffusion-weighted imaging using a unified pipeline. Performance is evaluated via patient-grouped five-fold cross-validation, employing strictly in-fold feature selection, probability calibration, and aligned out-of-fold predictions shared across all models. Furthermore, this study quantifies the incremental value of four routinely available clinical variables (age, PSA, PSA density, and prostate volume) and compares simple feature concatenation against more complex fusion strategies (dual-branch deep fusion and, for TabFM, dual-stream late fusion). Supported by comprehensive interpretability analyses—including feature-selection stability, cross-model feature importance, and spatial back-projection of radiomic descriptors—the primary objectives are to determine whether deep tabular models or a pretrained tabular foundation model confer a genuine discrimination advantage over strong tree-ensemble baselines, to test whether any such advantage depends on the foundation model's pretraining, and to assess exactly how much routinely collected clinical information adds to a calibrated, radiomics-based csPCa risk model.

\section{Materials and Methods}
\label{sec:methods}

\begin{figure}[t]
  \centering
  \safeinclude[width=\textwidth]{Methodology.png}
  \caption{\textbf{Overview of the study pipeline.} Biparametric MRI studies
    (T2-weighted, ADC, high b-value DWI) and whole-gland segmentation masks are
    processed through modality-specific extraction pipelines to build a
    multimodal radiomic feature table. A predefined patient-grouped five-fold
    cross-validation split is shared by all experimental conditions.}
  \label{fig:pipeline}
\end{figure}

This study benchmarked classical machine learning (ML), deep tabular learning
(DL), and a pretrained tabular foundation model (TabFM) for binary
classification of clinically significant prostate cancer (csPCa) on
biparametric MRI, and quantified the value added by routine clinical variables.
Four conditions were evaluated (Figure~\ref{fig:pipeline}):
radiomics-only, clinical-only, and two strategies that combine both sources
(feature concatenation and dual-branch deep fusion). All conditions reused a
single predefined cross-validation partition and, where radiomic features were
included, a locked fold-specific feature-selection plan, so that performance
differences reflected only the added information and the fusion strategy rather
than the evaluation protocol. Reporting follows the TRIPOD+AI
statement~\citep{collins2024tripodai} and the radiomics quality
recommendations of~\citet{lambin2017radiomics}.

%% ---------------------------------------------------------------
\subsection{Study cohort and outcome definition}
\label{sec:cohort}

Biparametric MRI studies were drawn from the Prostate Imaging: Cancer AI
(PI-CAI) dataset~\citep{saha2023artificial}, a large multicenter collection of
prostate MRI studies with histopathological outcomes and expert-derived
annotations. The outcome was clinically significant prostate cancer, defined by
the histopathological threshold of International Society of Urological Pathology
(ISUP) grade group~$\geq 2$~\citep{epstein2016grade}: grade groups 0--1 were
labeled non-csPCa and grade groups $\geq 2$ csPCa. The cohort comprised the full
PI-CAI Public Training and Development Dataset (1{,}500 studies from 1{,}476
patients), of which 425 (28.3\%) were csPCa-positive
(Table~\ref{tab:cohort}). The five outer cross-validation folds were taken
directly from the official PI-CAI case-level partition, which groups studies at
the patient level so that all studies from a given patient fall in a single
fold; the small excess of studies over patients reflects patients with more than
one study. These folds were not constructed to hold out acquisition centers,
scanner vendors, or imaging protocols. The contributing center and scanner model
were recovered for all 1{,}500 studies from the imaging provenance metadata,
identifying three centers---Radboud University Medical Center (RUMC, Nijmegen;
$n=800$), Prostaat Centrum Noord-Nederland (PCNN, Groningen; $n=350$), and
Ziekenhuisgroep Twente (ZGT, Hengelo; $n=350$)---and seven distinct scanner
models; PI-RADS scores were not recoverable at the case or lesion level from the
public metadata, marksheets, or delineation files. The cohort was heterogeneous
across centers: csPCa prevalence (31.1\%, 29.5\%, and 22.9\% for PCNN, RUMC, and
ZGT; $\chi^2$ $P=0.029$), patient age, serum PSA, and prostate volume all
differed significantly between centers (Kruskal--Wallis $P<10^{-5}$), whereas
PSA density did not ($P=0.48$). Because the official folds group patients but are
not center-disjoint, the primary five-fold evaluation should be interpreted as
multicenter internal validation with center-stratified subgroup performance
(Section~\ref{sec:res-center}); transportability to an unseen center was assessed
separately by a leave-one-center-out retraining analysis
(Section~\ref{sec:res-transport}). Radiomic features were extracted throughout
from a single, uniform automated whole-gland segmentation (Bosma22b) for every
study, so that the lesion-delineation provenance present in the dataset did not
enter the radiomic feature space.

\begin{table}[htbp]
\centering
\caption{Cohort characteristics. csPCa is defined as ISUP grade group $\geq 2$.
Continuous variables are median (interquartile range), compared with the
Mann--Whitney U test.}
\label{tab:cohort}
\begin{tabular}{lccccc}
\toprule
Variable & Overall & Non-csPCa & csPCa & Missing, n & $P$ \\
\midrule
Cases, n & 1500 & 1075 & 425 & 0 & \\
Patients, n & 1476 & 1057 & 425 & 0 & \\
Age, years & 66 (61--70) & 65 (60--69) & 68 (64--73) & 0 & $<0.001$ \\
PSA, ng/mL & 8.5 (5.9--13.0) & 7.8 (5.4--11.8) & 10.0 (6.8--17.0) & 40 & $<0.001$ \\
PSA density, ng/mL/cm$^3$ & 0.14 (0.09--0.22) & 0.12 (0.08--0.18) & 0.21 (0.13--0.36) & 451 & $<0.001$ \\
Prostate volume, cm$^3$ & 57 (40--80) & 61 (43--87) & 48 (35--64) & 27 & $<0.001$ \\
\midrule
ISUP grade group 0, n (\%) & 847 (56.5) & 847 & 0 & 0 & \\
ISUP grade group 1, n (\%) & 228 (15.2) & 228 & 0 & 0 & \\
ISUP grade group 2, n (\%) & 234 (15.6) & 0 & 234 & 0 & \\
ISUP grade group 3, n (\%) & 99 (6.6) & 0 & 99 & 0 & \\
ISUP grade group 4, n (\%) & 40 (2.7) & 0 & 40 & 0 & \\
ISUP grade group 5, n (\%) & 52 (3.5) & 0 & 52 & 0 & \\
\bottomrule
\end{tabular}
\end{table}

%% ---------------------------------------------------------------
\subsection{Image acquisition and radiomic feature extraction}
\label{sec:extraction}

Three image series were used: T2-weighted imaging (T2w), the apparent diffusion
coefficient (ADC) map derived from multi-b-value diffusion-weighted MRI, and
high b-value diffusion-weighted imaging (DWI). T2w
provides high-resolution zonal anatomy and is the primary sequence for lesion
localization under PI-RADS v2.1~\citep{turkbey2019pirads}; the ADC map is a
quantitative measure of tissue diffusivity, reduced in hypercellular tumors; and
high b-value DWI is sensitive to the restricted diffusion characteristic of
csPCa. Radiomic features were computed at the study level from the whole-gland
region of interest, using whole-gland masks provided with the dataset, framing
the task as study-level binary classification.

Features were extracted with PyRadiomics~v3~\citep{vanGriethuysen2017} following
the Image Biomarker Standardization Initiative~\citep{zwanenburg2020image}.
Extraction settings were tailored per sequence (Table~\ref{tab:extraction}) but
used an identical set of image filters and texture feature families across
sequences, so that each modality contributed comparable descriptors; ADC
normalization was disabled to preserve the quantitative interpretation of
diffusion coefficients. Modality-specific tables were merged into a single
multimodal feature table on a common study identifier; when clinical variables were included, age,
prostate-specific antigen (PSA), PSA density, and prostate volume were appended
to this table.

\begin{table}[htbp]
\centering
\caption{Radiomic feature extraction settings per MRI sequence (PyRadiomics).}
\label{tab:extraction}
\begin{tabular}{lccc}
\toprule
Setting & T2w & ADC & DWI \\
\midrule
Intensity normalization & Yes (scale 100) & No & Yes (scale 100) \\
Post-normalization shift & $+300$ & --- & $+300$ \\
In-plane resampling (B-spline) & 0.42~mm & 1.85~mm & 1.85~mm \\
Forced 2D (axial) extraction & Yes & Yes & Yes \\
Discretization bin width & 11.93 & 61.27 & 37.05 \\
\midrule
\multicolumn{4}{@{}l@{}}{\emph{Applied identically to all three sequences:}}\\
\multicolumn{4}{@{}l@{}}{\quad Image types: original; Laplacian of Gaussian ($\sigma = 0.6, 2, 3$~mm); level-2 wavelet;}\\
\multicolumn{4}{@{}l@{}}{\quad\quad square; square root; logarithm; exponential}\\
\multicolumn{4}{@{}l@{}}{\quad Texture families: first-order, GLCM, GLRLM, GLSZM, GLDM, NGTDM}\\
\multicolumn{4}{@{}l@{}}{\quad Shape features: extracted once, from T2w only}\\
\bottomrule
\end{tabular}
\end{table}

%% ---------------------------------------------------------------
\subsection{Feature selection and cross-validation}
\label{sec:selection}

Performance was estimated with the official PI-CAI patient-grouped five-fold
cross-validation, shared across every condition and across the ML and DL
pipelines. For each fold, all preprocessing, feature selection, hyperparameter
optimization, probability calibration, and (for deep models) early-stopping
decisions were performed using the training partition only, and performance was
evaluated on the held-out validation partition, which was never observed during
fitting or model selection. Out-of-fold predictions were accumulated across all
five folds to form a complete set of study-level calibrated probabilities for
aggregate analysis.

Feature selection was performed independently within each training fold to
prevent leakage. Each candidate feature was tested for association with the
csPCa label using a distribution-appropriate univariate test (Welch's $t$-test
for approximately normal features, otherwise the Mann--Whitney U
test~\citep{mann1947test}), with Benjamini--Hochberg false-discovery-rate
control at $\alpha = 0.05$~\citep{benjamini1995controlling}. Surviving features
were pruned for redundancy by removing, within each pair exceeding a Pearson
correlation of $0.95$, the feature with the lower univariate discrimination, and
the remaining features were capped at a data-adaptive limit derived from the
fold sample size and minority-class count (minimum 30, maximum 100 features).
The resulting fold-specific feature set was shared identically by the classical
and deep models, so observed differences were not attributable to different
feature subsets. Clinical variables, when present, bypassed this
filter and were retained unconditionally, preserving their interpretability as
established urological signals. Missing PSA density was first completed
deterministically as the ratio of PSA to prostate volume whenever both source
values were available; any remaining missing clinical values were then imputed by
the training-fold median, and all clinical variables were standardized using
training-fold statistics, with both the derivation and the imputation/scaling
statistics transferred unchanged to the validation fold. This deterministic
derivation recovered 399 of the 451 studies with missing PSA density, leaving only
52 studies (3.5\%) with a training-fold-median-imputed value; to confirm that the
imputation did not drive any conclusion, all clinical-containing conditions were
re-run on a strict source-complete cohort restricted to the 1{,}049 studies with a
non-derived, originally recorded PSA density (Section~\ref{sec:res-transport}).

%% ---------------------------------------------------------------
\subsection{Predictive models}
\label{sec:models}

Three classical ensemble algorithms served as the radiomics-only ML baselines:
Random Forest~\citep{breiman2001random}, Gradient
Boosting~\citep{friedman2001greedy}, and LightGBM~\citep{ke2017lightgbm}. These
were chosen from a preliminary screening of twelve classifiers (including
support vector machines, regularized logistic regression, discriminant analysis,
$k$-nearest neighbors, naive Bayes, decision tree, AdaBoost, and Extra Trees)
evaluated by repeated patient-grouped cross-validation (five folds, ten
repetitions); the three highest-ranked models by pooled out-of-fold AUROC were
promoted to the final benchmark. This screening used the same 1{,}500-study
cohort as the final five-fold evaluation rather than an independent development
sample. Consequently, comparisons involving the retained algorithms or
data-selected condition winners are treated as exploratory, and their confidence
intervals are conditional on the selected model set. For each fold,
hyperparameters were tuned by
randomized search with inner cross-validation, class imbalance was addressed by
inverse-frequency weighting, and predicted probabilities were calibrated by
Platt scaling~\citep{platt1999probabilistic} fit on held-out training-fold
predictions.

A pretrained tabular foundation model, TabFM~\citep{kong2026tabfm}
(weights \texttt{google/tabfm-1.0.0-pytorch}), was added as a fourth model
family alongside the classical ensembles and the deep architectures. Building on
the in-context-learning paradigm introduced by
TabPFN~\citep{hollmann2025tabpfn}, TabFM
performs inference by in-context learning over the training examples of each
fold rather than by dataset-specific gradient training, drawing on parameters
pretrained across a large corpus of synthetic tabular tasks; it was
evaluated on exactly the same fold-specific, leakage-safe feature-selection
output as the classical and deep models, so that any performance difference
reflects the model itself rather than a different feature representation.
TabFM was applied to the radiomics-only, clinical-only, and concatenation
conditions using the same outer folds and outer-training-derived Youden
threshold as the other model families. For the combined-modality condition,
TabFM was additionally evaluated with a dual-stream late-fusion strategy that
differs architecturally from the shared-dense-classifier dual-branch design
used for the deep models: independent TabFM instances were fit in-context on
the radiomics-only and clinical-only feature sets of each outer-training fold,
and their predicted probabilities on an inner-validation split of that same
fold (never including outer-fold cases) were combined by a logistic regression
fit on five derived inputs (the two probabilities, their logits, and their
absolute difference); this small fusion model and its Youden threshold, fit
only on the inner-validation split, were then applied unchanged to the two
TabFM streams' outer-fold test predictions. Because this dual-stream late
fusion is architecturally distinct from the shared-dense-classifier dual-branch
strategy used for the classical and deep models, the two ``dual'' conditions
are reported alongside each other but should not be read as a controlled
architectural comparison of fusion strategy. To determine whether TabFM's performance depended
on its pretrained parameters rather than merely its architecture and
in-context-learning procedure, a random-initialization control was evaluated in
the radiomics-only condition, in which the pretrained weights were replaced by
random initialization while every other component of the pipeline, including
the input features and folds, was left unchanged.

The deep architectures were chosen to probe two distinct and complementary
inductive biases for tabular data while keeping a controlled comparison against
the classical ensembles. Radiomic feature vectors are moderate-dimensional,
strongly inter-correlated, and informative chiefly through higher-order
interactions among intensity, texture, and shape descriptors; a plain multilayer
perceptron captures such interactions only implicitly and tends to overfit at
the present sample size. We therefore selected, first, a tabular
Transformer, whose self-attention reweights each feature in the context of all
others and thereby models feature--feature interactions explicitly; attention-based
models are the strongest contemporary deep competitors to gradient-boosted trees
on tabular benchmarks~\citep{gorishniy2021revisiting,vaswani2017attention} and
are thus the natural deep counterpart to the classical baselines. Second, a
Capsule Network, which replaces scalar activations with vector-valued capsules
refined by dynamic routing~\citep{sabour2017dynamic}; by encoding part--whole
relationships and instantiation parameters rather than mere feature presence, it
offers a different representational prior that may group related radiomic
descriptors and improve robustness in the low-sample regime typical of
radiomics. Third, a Transformer--CapsNet hybrid that composes the two
biases---attention-based interaction modeling followed by routing-based
agreement---to test whether they are complementary. Together the three span the
leading paradigms for deep tabular learning while sharing the feature space and
outer folds of the classical models.

These three deep tabular architectures were evaluated on the same selected
features. The tabular
Transformer~\citep{gorishniy2021revisiting,vaswani2017attention} maps the input
through GELU-activated dense layers~\citep{hendrycks2016gaussian} with batch
normalization~\citep{ioffe2015batch} into a sequence of learned tokens, applies
two multi-head self-attention blocks with a learnable positional embedding, and
aggregates the sequence by attention pooling before a sigmoid output; it is
trained with focal loss~\citep{lin2017focal} using the Adam
optimizer~\citep{kingma2015adam} with decoupled weight decay and cosine decay
with warm restarts~\citep{loshchilov2017sgdr}. The
CapsNet~\citep{sabour2017dynamic} encodes features into primary capsules and
refines two class capsules through dynamic routing, trained with margin loss,
and the Transformer--CapsNet hybrid feeds the Transformer token sequence into the
capsule routing stage. For each deep model and outer fold, an internal 80:20
split of the outer-training partition was first used for early stopping, Platt
scaling, and operating-threshold selection. A freshly initialized model was then
trained on the complete outer-training partition for the selected number of
epochs. The calibration map and threshold estimated without outer-fold labels
were applied to this refitted model's held-out predictions. Thus, the final
classical and deep estimators both used the complete outer-training partition,
although their inner hyperparameter- and epoch-selection procedures differed.

Routinely available clinical variables (age, PSA, PSA density, prostate volume)
were evaluated in three additional conditions on the same folds and protocols.
In the clinical-only condition, models were trained on the four variables alone.
In the concatenation condition, the variables were appended to the fold-specific
radiomic subset and passed to the same single-stream ML and DL models without
architectural change. In the dual-branch condition, clinical and radiomic
inputs were processed by two independent branches (no shared weights before
fusion) and combined by a shared dense classifier, implemented for each of the
three deep architectures; this enforces modality-specific latent representations
prior to fusion, in contrast to concatenation, which mixes the two sources from
the first layer.

%% ---------------------------------------------------------------
\subsection{Calibration, evaluation, and statistical analysis}
\label{sec:eval}

All models produced calibrated probabilities of csPCa via Platt
scaling~\citep{platt1999probabilistic}, with calibration assessed by reliability
diagrams and the Brier score~\citep{brier1950verification}. For classical
models, calibration and threshold selection used grouped inner-cross-validation
out-of-fold predictions from the outer-training partition; for deep models,
they used the internal validation subset of the outer-training partition. A
fold-specific operating point was selected on these outer-training predictions
using the Youden index~\citep{youden1950index},
\begin{equation}
  \hat{\tau} = \arg\max_{\tau}\,[\mathrm{Se}(\tau) + \mathrm{Sp}(\tau) - 1],
  \label{eq:youden}
\end{equation}
and was then applied unchanged to the corresponding held-out outer-fold
predictions. Thus, outer-fold labels were not used for calibration or threshold
selection. All threshold-dependent metrics are reported at this operating
point, the fixed 0.5 threshold serving only as a reference. Discrimination was
summarized by AUROC and AUPRC; AUPRC was treated as a co-primary metric given
the 28.3\% class prevalence. Calibration was summarized by the Brier score, and
clinical utility by decision-curve analysis~\citep{vickers2006decision}.
Sensitivity, specificity, balanced accuracy, predictive values, F1 score, and
the Matthews correlation coefficient~\citep{matthews1975comparison} were also
computed.

Each metric was computed once on the pooled out-of-fold predictions, and 95\%
confidence intervals were obtained by a patient-level cluster bootstrap (5000
resamples of whole patients), which was also used for pairwise model comparisons
of the metric difference (two-sided bootstrap $P$); the DeLong paired test was
reported as a reference for AUROC differences. Six contrasts were specified
\emph{a priori} and examined for both AUROC and AUPRC. Two were designated
primary because they directly address the two study aims: the radiomics-only
Random Forest--Transformer comparison, which tests whether deep tabular learning
confers a discrimination advantage over the strongest classical ensemble; and
the concatenated-versus-radiomics-only Random Forest comparison, which tests the
incremental value of routine clinical variables. The remaining four contrasts
(dual-branch versus radiomics-only, dual-branch versus concatenation,
radiomics-only versus clinical-only, and dual-branch versus clinical-only) were
secondary. This primary/secondary designation was analytical rather than
preregistered. Holm adjustment was applied across all six contrasts separately
for AUROC and AUPRC, and both raw and adjusted $P$ values are reported. The
seven contrasts involving TabFM were specified only after the foundation model
became available, later than the original six, and are therefore treated as a
separate Holm family adjusted among themselves rather than pooled with the
original set, so that the adjusted $P$ values of the pre-specified comparisons
are not altered retrospectively by the addition of a new model; this choice is
conservative within each family but does not control error across the two
families jointly, consistent with the overall exploratory framing. The
patient-level cluster bootstrap was the primary inferential tool because a small
number of patients contribute more than one study, which mildly violates the
independence assumption of standard tests; the DeLong test is reported only as a
familiar reference for AUROC differences. Because the preliminary screening of
twelve classical algorithms and the selection of the per-condition winning models
both reused this single cohort, all hypothesis tests are interpreted as
exploratory: the reported confidence intervals are conditional on the selected
model set and do not incorporate model-selection uncertainty, and confirmatory
inference would require nested model selection or an independent validation cohort.

To characterize robustness across acquisition sites, discrimination was also
computed within each center on the pooled out-of-fold predictions, with
patient-level cluster-bootstrap 95\% confidence intervals. Because the three
centers are disjoint sets of patients, the difference in AUROC between the
lowest-performing center and the remaining centers was assessed with an unpaired
patient-level bootstrap (5000 resamples drawn independently within each group),
reported as the mean difference, 95\% confidence interval, and two-sided
bootstrap $P$ value. Cohort heterogeneity across centers was summarized with the
$\chi^2$ test for categorical variables and the Kruskal--Wallis test for
continuous variables.

Because the official patient-grouped folds mix acquisition centers and therefore
report per-center performance of a center-mixed--trained model rather than
transportability to an unseen center, discrimination was additionally estimated
under a leave-one-center-out (LOCO) protocol. In each of the three LOCO folds the
model was trained on two centers and tested on the entirely held-out third center
(train RUMC$+$ZGT / test PCNN; train PCNN$+$ZGT / test RUMC; train PCNN$+$RUMC /
test ZGT); all preprocessing, fold-specific feature selection, calibration, and
threshold selection were confined to the two training centers, and for TabFM the
in-context background rows were drawn only from the training centers, so the
held-out center never entered any fitting or context set. Out-of-fold predictions
were pooled across the three held-out centers to yield a single transportability
estimate per condition (radiomics-only, clinical-only, and radiomics$+$clinical) for
the classical, deep, and TabFM families. Two further sensitivity analyses probed
the robustness of the primary five-fold results: (i)~a source-complete PSA-density
analysis restricted to the 1{,}049 studies with an originally recorded PSA density
(versus the full cohort with deterministic derivation and residual median
imputation), and (ii)~a calibration-method comparison replacing Platt scaling with
isotonic regression for the classical radiomics-only and concatenation models.
Feature-selection stability was summarized by the per-fold selection frequency and
the pairwise Jaccard overlap of the five fold-specific feature sets, and clinical
utility was expressed, from the decision-curve analysis, as the net benefit and the
number of biopsies avoided per 100 patients relative to a biopsy-all strategy at
each model's operating point.

%% ---------------------------------------------------------------
\subsection{Interpretability and feature-stability analysis}
\label{sec:interpret-methods}

Because a central aim was to characterize \emph{which} imaging information drives
the prediction and how reproducible it is, interpretability was assessed at three
complementary levels. At the level of the selected biomarker signature, the
stability of in-fold feature selection was quantified by the number and source
sequence of features chosen per fold and by their reproducibility across folds,
summarized by selection frequency and the mean pairwise Jaccard overlap of the
five fold-specific feature sets. At the level of predictive contribution, global
feature importance was computed for every model: model-native importance
(impurity- or gain-based for the tree ensembles), SHAP
values~\citep{lundberg2017unified} for the classical models, and integrated
gradients~\citep{sundararajan2017axiomatic} for the deep models, together with
permutation importance~\citep{fisher2019all} as a model-agnostic common currency
computed identically across families as the mean decrease in out-of-fold AUROC
when a feature is randomly permuted. For TabFM, whose in-context inference
exposes no differentiable weights, model-agnostic permutation
SHAP~\citep{lundberg2017unified} was computed on a held-out subsample of each
fold (25 validation studies per fold, with a 50-study background set), yielding
per-study signed attributions in the radiomics-only, clinical-only, and
concatenation conditions. The concordance of the leading features
across model families was summarized by the pairwise top-15 Jaccard overlap and
by the features shared by all models. Finally, at the spatial level, the
individual high-importance features were back-projected to the image as
voxel-wise radiomic maps, computed within the gland and overlaid on the
corresponding ADC or DWI series together with the prostate-gland and
csPCa-lesion delineations. For spatial comparison, each displayed case also
included the original DWI slice with only the gland and lesion contours. Feature
maps were rendered on a relative low-to-high scale for each feature panel, so
that the spatial pattern of each descriptor could be inspected without implying
that raw radiomic units are comparable across feature families.

%% ---------------------------------------------------------------
\section{Results}
\label{sec:results}

The analysis cohort comprised 1{,}500 studies from 1{,}476 patients, of which
425 (28.3\%) were csPCa-positive; the groups differed in all four clinical
variables (Table~\ref{tab:cohort}). All discrimination and calibration metrics
below were computed on aligned pooled out-of-fold predictions with 95\%
patient-level cluster-bootstrap confidence intervals, and threshold-dependent
metrics use each model's per-fold outer-training-derived Youden operating point
rather than an arbitrary 0.5 cut, which at the cohort prevalence would
substantially understate sensitivity.

\subsection{Radiomics-only classification: classical versus deep models}
\label{sec:res-radiomics}

On the radiomics-only feature space, TabFM ranked first on both discrimination
metrics (AUROC 0.828, 95\% CI 0.805--0.850; AUPRC 0.650, 0.600--0.699), above
the classical Random Forest (0.814; 0.605) and all three deep models (AUROC
0.794--0.803), with the lowest Brier score (0.145); the full ranking with
confidence intervals is shown in Figure~\ref{fig:radiomics_forest} and
per-model operating-point metrics in Table~S2 (Supplementary Material). The
confidence intervals overlapped substantially across families: in the specified
TabFM--Random-Forest comparison the AUROC difference was not significant after
multiplicity adjustment ($\Delta$AUROC $+0.014$; Holm-adjusted $P=0.084$),
although AUPRC was ($\Delta=+0.045$; Holm-adjusted $P=0.022$;
Table~\ref{tab:significance}), and the originally specified Random
Forest--Transformer comparison was likewise non-significant ($\Delta$AUROC
$+0.014$; Holm-adjusted $P=0.272$). At each model's outer-training-derived
Youden threshold (median 0.18--0.26) the models operated at a clinically
balanced point, with high negative predictive value (0.86--0.90) but only
moderate positive predictive value (0.49--0.54) at the 28.3\% disease
prevalence---so the models are more informative for ruling out than for
confirming csPCa. A random-initialization
control, in which TabFM's pretrained parameters were replaced with random
weights while every other component of the pipeline was left unchanged,
collapsed to near-chance discrimination (AUROC 0.522, AUPRC 0.294) and, at the
0.5 threshold, defaulted to predicting every study as csPCa-positive
(sensitivity 1.000, specificity 0.000). This control bounds rather than fully
isolates the effect of pretraining---an in-context model with random weights is
not expected to be competitive---but, because the architecture and
in-context-learning procedure were otherwise unchanged, it confirms that
TabFM's radiomics-only performance depends on its pretrained parameters rather
than on its architecture or in-context mechanism alone.

\begin{figure}[!htbp]
\centering
\safeinclude[width=\textwidth]{fig_radiomics_forest.png}
\caption{\textbf{Radiomics-only discrimination across model families.}
Pooled out-of-fold (A) AUROC and (B) AUPRC with patient-level cluster-bootstrap
95\% confidence intervals for the seven radiomics-only models, ordered by AUROC.
The pretrained tabular foundation model (TabFM, orange) ranks first on both
metrics, but its confidence intervals overlap those of the classical and deep
baselines.}
\label{fig:radiomics_forest}
\end{figure}


\subsection{Added value of clinical variables}
\label{sec:res-clinical}

The four clinical variables alone were markedly less discriminative than
radiomics (best clinical-only AUROC 0.757, 95\% CI 0.730--0.784, achieved by
TabFM; radiomics-only Random Forest versus clinical-only Transformer--CapsNet
$\Delta$AUROC $+0.060$, Holm-adjusted $P<0.001$; TabFM's own clinical-only
result was statistically indistinguishable from Transformer--CapsNet,
$\Delta$AUROC $+0.003$, Holm-adjusted $P=0.921$; Table~\ref{tab:hybrid}).
Combining the two sources improved discrimination over radiomics alone,
producing a consistent upward shift in both the ROC and precision--recall
curves (Figure~\ref{fig:discrimination}). Feature concatenation passed to TabFM
reached the best overall AUROC of 0.844 (0.822--0.865), a statistically
significant gain over the previous best concatenation model (Random Forest,
$\Delta$AUROC $+0.016$, 95\% CI 0.005--0.026, raw $P=0.003$, Holm-adjusted
$P=0.017$; $\Delta$AUPRC $+0.032$, Holm-adjusted $P=0.087$), over the best
dual-branch model (Transformer, $\Delta$AUROC $+0.018$, Holm-adjusted
$P=0.008$), and over the radiomics-only Random Forest ($\Delta$AUROC $+0.030$,
Holm-adjusted $P<0.001$; $\Delta$AUPRC $+0.061$, Holm-adjusted $P=0.004$; these
TabFM contrasts form a separate family from the six originally specified
contrasts and were Holm-adjusted among themselves; Table~\ref{tab:significance}).
Radiomics-only TabFM was itself a significant improvement over the
radiomics-only Random Forest in AUPRC, as reported above, and the concatenated
Random Forest remained a significant improvement over the radiomics-only Random
Forest ($\Delta$AUROC $+0.014$, 95\% CI 0.007--0.021, raw $P<0.001$,
Holm-adjusted $P<0.001$; $\Delta$AUPRC $+0.029$, raw $P=0.004$, Holm-adjusted
$P=0.022$). Among the classical and deep architectures, the best dual-branch
model (Transformer, shared-dense-classifier fusion) reached AUROC 0.826
(0.802--0.848) and AUPRC 0.638 (0.589--0.692), and neither its difference from
radiomics-only Random Forest nor from concatenated Random Forest remained
significant after Holm adjustment (AUROC $P=0.266$ and $P=0.736$,
respectively), so the added complexity of dual-branch fusion conferred no
demonstrated benefit over simple concatenation among these architectures. When
TabFM was instead combined with clinical variables through the dual-stream
late-fusion strategy described in Methods, it reached AUROC 0.844
(0.822--0.865) and AUPRC 0.674 (0.626--0.726), a significant improvement over
the dual Transformer ($\Delta$AUROC $+0.018$, Holm-adjusted $P=0.004$;
$\Delta$AUPRC $+0.036$, Holm-adjusted $P=0.022$) that was numerically
indistinguishable from the concatenated TabFM ($\Delta$AUROC $+0.0001$,
Holm-adjusted $P=1.00$; $\Delta$AUPRC $+0.008$, Holm-adjusted $P=1.00$;
Table~\ref{tab:significance}), so, unlike the classical and deep architectures,
TabFM showed no advantage from either fusion strategy over the other: its
concatenation and late-fusion results were essentially the same model in
performance terms. Because the two ``dual'' strategies differ architecturally
(shared-dense-classifier fusion for the classical and deep models versus
logistic late fusion of two independently in-context-trained TabFM streams),
this equivalence should not be read as evidence that TabFM's fusion mechanism
is inherently more effective than dual-branch fusion, only that neither
fusion strategy improved on simple concatenation within its own model family.
At their selected operating points the combined models behaved similarly to the
radiomics-only models, with balanced accuracy of 0.75--0.76, except that the
dual-fusion TabFM's Youden threshold (median 0.53) shifted its operating point
toward higher specificity (0.761) than sensitivity (0.734), unlike the other
combined models. The concatenated TabFM had sensitivity 0.767 and specificity
0.745, comparable to the concatenated Random Forest (sensitivity 0.774,
specificity 0.729) and the dual Transformer (sensitivity 0.746, specificity
0.741). Reliability curves, Brier scores, and decision curves are shown in
Figure~\ref{fig:calibration_dca}: the concatenated TabFM had the lowest Brier
score among all combined models (0.139), whereas the dual-fusion TabFM's Brier
score (0.167) was the highest among the combined models, reflecting weaker
calibration of the small inner-validation-fit fusion model rather than weaker
discrimination. The improvement from adding clinical variables was modest
across model families and clearest when routed through TabFM; confidence
intervals for the operating-point statistics overlapped.

\begin{table}[!htbp]
\centering
\caption{Best model in each condition on pooled out-of-fold predictions, with
threshold-dependent metrics at the outer-training-derived Youden operating point.
Threshold-free metrics use the calibrated probabilities; threshold-dependent
metrics use the per-fold outer-training-derived Youden operating point, with
patient-level cluster-bootstrap 95\% confidence intervals for AUROC, AUPRC,
sensitivity, and specificity (lower Brier is better).
For the dual-branch/dual-stream condition, TabFM's logistic late-fusion result
is architecturally distinct from the shared-dense-classifier dual-branch design
used for the classical and deep models (see Methods).}
\label{tab:hybrid}
\resizebox{\textwidth}{!}{
\begin{tabular}{llccccccc}
\toprule
Condition & Best model & AUROC & AUPRC & Brier & Sensitivity & Specificity & Bal.\ acc. & NPV \\
\midrule
Clinical-only & TabFM & 0.757 (0.730--0.784) & 0.533 (0.487--0.588) & 0.169 & 0.633 (0.587--0.679) & 0.753 (0.727--0.780) & 0.693 & 0.839 \\
Radiomics-only & TabFM & 0.828 (0.805--0.850) & 0.650 (0.600--0.699) & 0.145 & 0.812 (0.774--0.848) & 0.691 (0.664--0.718) & 0.751 & 0.903 \\
Radiomics$+$clinical (concat) & TabFM & 0.844 (0.822--0.865) & 0.666 (0.618--0.719) & 0.139 & 0.767 (0.726--0.806) & 0.745 (0.720--0.771) & 0.756 & 0.890 \\
Radiomics$+$clinical (dual) & TabFM (late fusion) & 0.844 (0.822--0.865) & 0.674 (0.626--0.726) & 0.167 & 0.734 (0.691--0.776) & 0.761 (0.735--0.787) & 0.748 & 0.879 \\
\bottomrule
\end{tabular}
}
\end{table}

\begin{figure}[!htbp]
\centering
\safeinclude[width=\textwidth]{fig_discrimination.png}
\caption{\textbf{Discrimination of the best model in each condition} on aligned
out-of-fold predictions. (A) Receiver operating characteristic curves with pooled
AUROC. (B) Precision--recall curves with pooled AUPRC; the dotted line marks the
csPCa prevalence (0.283). Combining radiomics with clinical variables shifts both
curves upward relative to radiomics alone, while clinical variables alone are
clearly inferior.}
\label{fig:discrimination}
\end{figure}

\begin{figure}[!htbp]
\centering
\safeinclude[width=\textwidth]{fig_calibration_dca.png}
\caption{\textbf{Calibration and clinical utility of the best model in each
condition.} (A) Reliability curves (10 quantile bins); the diagonal denotes
perfect calibration. (B) Decision-curve analysis showing net benefit versus
threshold probability, with treat-all and treat-none reference strategies.}
\label{fig:calibration_dca}
\end{figure}

\begin{table}[!htbp]
\centering
\caption{Exploratory pairwise model comparisons (A versus B) on aligned
out-of-fold predictions. $\Delta$ is A$-$B; 95\% CIs and raw bootstrap $P$
values are from a patient-level cluster bootstrap. Holm adjustment was applied
across the six originally specified contrasts separately for AUROC and AUPRC.
The seven TabFM contrasts (lower block) were added after the foundation model
became available and form a separate Holm family, adjusted among themselves and
not pooled with the original six. DeLong $P$ is the paired AUROC reference
test. Confidence intervals are unadjusted and conditional on the screened
model set.}
\label{tab:significance}
\resizebox{\textwidth}{!}{
\begin{tabular}{lcccc}
\toprule
Contrast (A vs.\ B) & $\Delta$AUROC [95\% CI] & Bootstrap $P$ raw (Holm) & DeLong $P$ raw (Holm) & $\Delta$AUPRC [95\% CI], $P$ raw (Holm) \\
\midrule
Concat (RF) vs.\ radiomics-only (RF) & $+0.014$ [0.007, 0.021] & $<0.001$ ($<0.001$) & $<0.001$ ($<0.001$) & $+0.029$ [0.010, 0.050], 0.004 (0.022) \\
Dual (Transformer) vs.\ radiomics-only (RF) & $+0.012$ [$-0.004$, 0.027] & 0.133 (0.266) & 0.131 (0.261) & $+0.033$ [$-0.006$, 0.072], 0.092 (0.277) \\
Dual (Transformer) vs.\ concat (RF) & $-0.002$ [$-0.015$, 0.011] & 0.736 (0.736) & 0.727 (0.727) & $+0.004$ [$-0.026$, 0.034], 0.748 (1.000) \\
Radiomics-only RF vs.\ Transformer & $+0.014$ [$-0.001$, 0.029] & 0.068 (0.204) & 0.071 (0.213) & $-0.004$ [$-0.037$, 0.030], 0.829 (1.000) \\
Radiomics-only (RF) vs.\ clinical-only (Transformer--CapsNet) & $+0.060$ [0.029, 0.091] & $<0.001$ ($<0.001$) & $<0.001$ ($<0.001$) & $+0.057$ [0.002, 0.114], 0.042 (0.170) \\
Dual (Transformer) vs.\ clinical-only (Transformer--CapsNet) & $+0.072$ [0.049, 0.094] & $<0.001$ ($<0.001$) & $<0.001$ ($<0.001$) & $+0.090$ [0.052, 0.130], $<0.001$ ($<0.001$) \\
\midrule
\multicolumn{5}{@{}l}{\emph{TabFM contrasts (separate Holm family)}} \\
Concat (TabFM) vs.\ concat (RF) & $+0.016$ [0.005, 0.026] & 0.003 (0.020) & 0.004 (0.027) & $+0.032$ [0.007, 0.058], 0.012 (0.099) \\
Concat (TabFM) vs.\ dual (Transformer) & $+0.018$ [0.007, 0.029] & 0.001 (0.010) & 0.001 (0.010) & $+0.028$ [0.004, 0.051], 0.022 (0.157) \\
Concat (TabFM) vs.\ radiomics-only (RF) & $+0.030$ [0.016, 0.043] & $<0.001$ ($<0.001$) & $<0.001$ ($<0.001$) & $+0.061$ [0.027, 0.095], $<0.001$ (0.005) \\
Radiomics-only (TabFM) vs.\ radiomics-only (RF) & $+0.014$ [0.003, 0.025] & 0.014 (0.084) & 0.012 (0.073) & $+0.045$ [0.015, 0.072], 0.002 (0.022) \\
Clinical-only (TabFM) vs.\ clinical-only (Transformer--CapsNet) & $+0.003$ [$-0.004$, 0.010] & 0.460 (1.000) & 0.462 (1.000) & $-0.015$ [$-0.029$, 0.003], 0.100 (0.462) \\
Dual-fusion (TabFM) vs.\ dual (Transformer) & $+0.018$ [0.007, 0.029] & $<0.001$ (0.004) & $<0.001$ (0.008) & $+0.036$ [0.012, 0.060], 0.002 (0.022) \\
Dual-fusion (TabFM) vs.\ concat (TabFM) & $+0.000$ [$-0.007$, 0.008] & 0.954 (1.000) & 0.974 (1.000) & $+0.008$ [$-0.009$, 0.025], 0.353 (1.000) \\
\bottomrule
\end{tabular}
}
\end{table}

\subsection{Performance across acquisition centers}
\label{sec:res-center}

Stratifying the pooled out-of-fold predictions by contributing center revealed
substantial between-center heterogeneity that was consistent across model
families, including TabFM (Table~\ref{tab:center}). For the radiomics-only Random Forest,
discrimination was high and similar at RUMC (AUROC
0.835, 95\% CI 0.804--0.863) and ZGT (0.834, 0.784--0.882) but markedly lower at
PCNN (0.735, 0.678--0.788). The same pattern held for the combined models and was
statistically significant: PCNN underperformed the other two centers by
$0.102$ AUROC for the radiomics-only Random Forest (95\% CI 0.041--0.163,
$P=0.001$), by $0.130$ for the concatenation Random Forest
($P<0.001$), and by $0.153$ for the dual-branch Transformer ($P<0.001$). The
best-performing model family, TabFM, showed the same site-dependent pattern
rather than escaping it: PCNN again lagged RUMC and ZGT in every condition, with
radiomics-only TabFM reaching AUROC 0.849 (RUMC) and 0.847 (ZGT) but only 0.748
(PCNN, 95\% CI 0.694--0.799), and the concatenation and dual-fusion TabFM models
reaching 0.87 at RUMC and ZGT but 0.734 and 0.743 at PCNN
(Table~\ref{tab:center}). The per-center confidence intervals for TabFM at PCNN
did not overlap the corresponding RUMC estimates in the combined conditions,
indicating that the foundation model's aggregate advantage did not close the
between-center gap. Notably,
the incremental value of clinical variables was also center-dependent: adding
clinical variables improved discrimination at RUMC and ZGT but not at PCNN, where
the combined models were no better than---and for dual-branch fusion slightly
worse than---the radiomics-only model, a pattern that held for TabFM as well.
Because the cross-validation folds are not
center-disjoint, these estimates reflect per-center performance of a
center-mixed--trained model and not transportability to a wholly unseen center;
the lower performance at PCNN nonetheless signals that acquisition-site
characteristics materially affect this radiomic signature, independent of the
modelling approach, and that
center-stratified evaluation is essential before deployment.

\begin{table}[!htbp]
\centering
\caption{Discrimination by contributing center on pooled out-of-fold predictions
for the best classical/deep model in the radiomics-only, concatenation, and
dual-branch conditions, and for TabFM in the radiomics-only, concatenation, and
dual-stream late-fusion conditions. Values are point estimates with
patient-level cluster-bootstrap 95\% confidence intervals. RUMC, Radboud
University Medical Center; PCNN, Prostaat Centrum Noord-Nederland; ZGT,
Ziekenhuisgroep Twente.}
\label{tab:center}
\resizebox{\textwidth}{!}{
\begin{tabular}{llccc}
\toprule
Model & Center & $n$ (csPCa \%) & AUROC & AUPRC \\
\midrule
\multirow{4}{*}{Radiomics-only (Random Forest)}
 & All  & 1500 (28.3) & 0.814 (0.791--0.837) & 0.605 (0.554--0.659) \\
 & RUMC & 800 (29.5)  & 0.835 (0.804--0.863) & 0.672 (0.606--0.733) \\
 & ZGT  & 350 (22.9)  & 0.834 (0.784--0.882) & 0.581 (0.469--0.700) \\
 & PCNN & 350 (31.1)  & 0.735 (0.678--0.788) & 0.520 (0.429--0.624) \\
\midrule
\multirow{4}{*}{Radiomics$+$clinical concat (Random Forest)}
 & All  & 1500 (28.3) & 0.828 (0.806--0.850) & 0.634 (0.583--0.687) \\
 & RUMC & 800 (29.5)  & 0.850 (0.821--0.877) & 0.706 (0.642--0.765) \\
 & ZGT  & 350 (22.9)  & 0.863 (0.816--0.906) & 0.642 (0.528--0.751) \\
 & PCNN & 350 (31.1)  & 0.725 (0.667--0.780) & 0.511 (0.422--0.614) \\
\midrule
\multirow{4}{*}{Radiomics$+$clinical dual (Transformer)}
 & All  & 1500 (28.3) & 0.826 (0.803--0.848) & 0.638 (0.589--0.691) \\
 & RUMC & 800 (29.5)  & 0.860 (0.830--0.886) & 0.740 (0.679--0.794) \\
 & ZGT  & 350 (22.9)  & 0.855 (0.807--0.899) & 0.635 (0.520--0.749) \\
 & PCNN & 350 (31.1)  & 0.706 (0.648--0.761) & 0.493 (0.405--0.598) \\
\midrule
\multirow{3}{*}{Radiomics-only (TabFM)}
 & RUMC & 800 (29.5)  & 0.849 (0.818--0.877) & 0.696 (0.634--0.759) \\
 & ZGT  & 350 (22.9)  & 0.847 (0.792--0.896) & 0.667 (0.569--0.772) \\
 & PCNN & 350 (31.1)  & 0.748 (0.694--0.799) & 0.557 (0.468--0.658) \\
\midrule
\multirow{3}{*}{Radiomics$+$clinical concat (TabFM)}
 & RUMC & 800 (29.5)  & 0.872 (0.844--0.898) & 0.739 (0.678--0.804) \\
 & ZGT  & 350 (22.9)  & 0.874 (0.826--0.915) & 0.700 (0.593--0.803) \\
 & PCNN & 350 (31.1)  & 0.734 (0.677--0.785) & 0.519 (0.436--0.628) \\
\midrule
\multirow{3}{*}{Radiomics$+$clinical late fusion (TabFM)}
 & RUMC & 800 (29.5)  & 0.874 (0.847--0.901) & 0.768 (0.714--0.821) \\
 & ZGT  & 350 (22.9)  & 0.870 (0.822--0.915) & 0.692 (0.586--0.788) \\
 & PCNN & 350 (31.1)  & 0.743 (0.688--0.794) & 0.530 (0.443--0.636) \\
\bottomrule
\end{tabular}
}
\end{table}

\subsection{Cross-center transportability and sensitivity analyses}
\label{sec:res-transport}

The center-stratified results above report a center-mixed--trained model evaluated
per center; to test genuine transportability we retrained every condition under a
leave-one-center-out (LOCO) protocol and pooled the predictions over the three
held-out centers (Table~\ref{tab:loco}). Discrimination fell substantially for every
model family relative to the center-mixed five-fold estimates. In the radiomics-only
setting the pooled LOCO AUROC dropped to 0.702 for TabFM (from 0.828), 0.695 for the
best classical model (LightGBM; Random Forest 0.676, from 0.814), and to a similar
range for the deep models; adding clinical variables recovered part of the loss
(concatenation pooled LOCO AUROC 0.753 for TabFM and 0.728 for the best classical
model), and clinical variables alone transported at 0.70--0.72. TabFM retained a
small numerical lead in every LOCO condition, and its relative robustness was
greatest in the combined condition ($+0.025$--$0.045$ AUROC over the best classical
model), but no model approached its internal cross-validation performance. The
implied optimistic bias of the center-mixed folds---roughly $0.10$--$0.14$ AUROC in
the radiomics-only and combined conditions---is consistent in magnitude with the
between-center gap identified at PCNN, and indicates that the aggregate five-fold
numbers should be read as multicenter internal validation rather than as estimates of
performance at a new site.

The primary conclusions were otherwise robust to the two pre-planned sensitivity
analyses. Restricting the analysis to the 1{,}049 studies with an originally recorded
(non-derived) PSA density left the ordering and magnitude of results essentially
unchanged and, if anything, slightly higher than the full-cohort imputed analysis
(radiomics-only Random Forest AUROC 0.832, concatenation Random Forest 0.857,
concatenation TabFM 0.873), confirming that neither the derived nor the
median-imputed PSA-density values inflated the combined-model performance
(Table~\ref{tab:psad}). Replacing Platt scaling with isotonic regression did not
change discrimination or the ranking of the classical models in either the
radiomics-only (Random Forest AUROC 0.813 vs.\ 0.814) or the concatenation condition
(Random Forest 0.824 vs.\ 0.828), as expected for a monotonic recalibration
(Table~\ref{tab:psad}). Feature selection was stable across folds, with a mean
pairwise Jaccard overlap of 0.67 for the radiomics-only plan and 0.69 for the
radiomics$+$clinical plan (Table~S6). Finally, at each model's outer-training-derived
operating point the decision-curve analysis translated into a clinically meaningful
reduction in biopsies: relative to a biopsy-all strategy, the concatenation TabFM
avoided 37 biopsies per 100 patients at a net-benefit gain of $+0.135$ and a negative
predictive value of 0.90, and the radiomics-only TabFM avoided 32 per 100 at
NPV~0.91 (Table~S7).

\begin{table}[!htbp]
\centering
\caption{Cross-center transportability under leave-one-center-out (LOCO)
retraining. AUROC and AUPRC are pooled over the three held-out centers with
patient-level cluster-bootstrap 95\% confidence intervals for AUROC. The
right-hand column repeats the center-mixed five-fold estimate from the main
analysis for reference; the difference ($\Delta$) is the optimistic bias of the
center-mixed folds. Best classical model is the top-ranked ensemble in each
condition; Random Forest is shown as the primary classical reference.}
\label{tab:loco}
\resizebox{\textwidth}{!}{
\begin{tabular}{llccc}
\toprule
Condition & Model & LOCO AUROC (95\% CI) & LOCO AUPRC & Five-fold AUROC ($\Delta$) \\
\midrule
\multirow{3}{*}{Radiomics-only}
 & Random Forest        & 0.676 (0.647--0.706) & 0.443 & 0.814 ($-0.138$) \\
 & LightGBM (best ML)   & 0.695 (0.663--0.725) & 0.470 & 0.792 ($-0.097$) \\
 & TabFM                & 0.702 (0.674--0.729) & 0.461 & 0.828 ($-0.126$) \\
\midrule
\multirow{3}{*}{Clinical-only}
 & Random Forest        & 0.697 (0.666--0.729) & 0.457 & --- \\
 & Gradient Boosting (best ML) & 0.700 (0.668--0.731) & 0.465 & --- \\
 & TabFM                & 0.718 (0.690--0.746) & 0.480 & 0.757 ($-0.039$) \\
\midrule
\multirow{3}{*}{Radiomics$+$clinical (concat)}
 & Random Forest        & 0.708 (0.678--0.738) & 0.485 & 0.828 ($-0.120$) \\
 & Gradient Boosting (best ML) & 0.728 (0.699--0.756) & 0.501 & --- \\
 & TabFM                & 0.753 (0.726--0.779) & 0.536 & 0.844 ($-0.091$) \\
\bottomrule
\end{tabular}
}
\end{table}

\begin{table}[!htbp]
\centering
\caption{Sensitivity of the primary five-fold results to PSA-density handling and
to the calibration method. All values are pooled out-of-fold AUROC. ``Primary''
is the full cohort ($n=1500$) with deterministic PSA-density derivation, residual
training-fold median imputation, and Platt scaling. ``Source-complete'' restricts
to the $n=1049$ studies with an originally recorded PSA density. ``Isotonic''
replaces Platt scaling with isotonic regression on the full cohort and was
evaluated for the classical models only.}
\label{tab:psad}
\begin{tabular}{llccc}
\toprule
Condition & Model & Primary ($n{=}1500$) & Source-complete ($n{=}1049$) & Isotonic ($n{=}1500$) \\
\midrule
\multirow{2}{*}{Radiomics-only}
 & Random Forest & 0.814 & 0.832 & 0.813 \\
 & TabFM         & 0.828 & 0.851 & --- \\
\midrule
\multirow{2}{*}{Radiomics$+$clinical (concat)}
 & Random Forest & 0.828 & 0.857 & 0.824 \\
 & TabFM         & 0.844 & 0.873 & --- \\
\bottomrule
\end{tabular}
\end{table}

\subsection{Interpretability}
\label{sec:res-interpret}

In-fold feature selection was compact and highly reproducible. A mean of 67.6
features was retained per fold (SD 0.9; range 67--69), drawn from 99 distinct
features across the five folds; 41 of these were selected in every fold and 60 in
at least four, with a mean pairwise Jaccard overlap of 0.68
(Figure~\ref{fig:stability}A), indicating a stable radiomic signature rather than
fold-dependent selection. The signature was strongly diffusion-weighted: of the
99 features, 80 were derived from DWI, 12 from T2w, and 7 from ADC
(Figure~\ref{fig:stability}B). Although fewer in number, ADC descriptors were a
consistent component (three to six were retained in every fold) and were not
incidental to prediction, recurring among the most important variables of all six
models.

Feature-attribution analyses for the best model of each condition characterized
both the magnitude and the direction of the leading predictors
(Figure~\ref{fig:beeswarm}). In the radiomics-only Random Forest, the dominant
features were diffusion-weighted texture and first-order descriptors: DWI
gray-level co-occurrence informational measure of correlation (Imc1) and
correlation, neighbouring-gray-tone-difference strength, and the skewness and
kurtosis of the high b-value signal. The clinical-only model distributed its
importance across all four variables in clinically expected directions: higher
prostate-specific antigen, higher PSA density, and older age increased the
predicted probability of csPCa, whereas larger prostate volume decreased it.
Crucially, when the two information sources were combined the clinical variables
rose to the top of the ranking: PSA density was the single most important feature
in the concatenation Random Forest (ahead of every radiomic descriptor, with
higher values increasing csPCa probability) and all four clinical
variables---age, PSA density, prostate-specific antigen, and prostate
volume---appeared among the leading features of the dual-branch Transformer,
interleaved with diffusion- and T2-weighted texture, including DWI gray-level
co-occurrence (Imc1), DWI first-order skewness, and T2-weighted gray-level
non-uniformity and coarseness. This re-ranking offers a direct explanation for the modest
discrimination gain conferred by the clinical variables: rather than refining the
radiomic signal, they contributed largely independent, high-value information
dominated by PSA density.

Effect directions were concordant between the SHAP (classical) and
integrated-gradient (deep) analyses for the features they shared. Among the
radiomics-only models, agreement on the ranking of individual features across all
six architectures was nonetheless only moderate (mean pairwise top-15 Jaccard
0.27; Random Forest and the Transformer shared 5 of their 15 leading features;
Figure~\ref{fig:concordance}), reflecting the redundancy of the high-dimensional
radiomic space; despite this, one diffusion-weighted gray-level co-occurrence
feature (wavelet-filtered correlation) ranked among the fifteen most important
features of every model, a signature shared across classical and deep
architectures. Because TabFM performs inference by in-context learning rather
than by a differentiable forward pass over a fixed set of weights, it was
attributed with model-agnostic permutation SHAP computed on a held-out
subsample of each fold, in the radiomics-only, clinical-only, and concatenation
conditions (Supplementary Material, Figure~S1). Reassuringly, TabFM recovered the same
signature as the other model families using this independent attribution
method. In the radiomics-only condition its most influential features were the
ADC square-root first-order total energy and diffusion-weighted texture,
including the wavelet-filtered DWI gray-level co-occurrence correlation feature
that ranked among the top fifteen of every other model, together with DWI
gray-level non-uniformity and informational measure of correlation (Imc2) and
ADC gray-level non-uniformity. In the clinical-only condition PSA density was by
a wide margin the most important variable, followed by age and prostate volume,
in the expected directions (higher PSA density and age, and lower prostate
volume, increasing the predicted probability). Most tellingly, in the
concatenation condition PSA density and age rose to the top two positions ahead
of every radiomic descriptor, exactly reproducing the re-ranking observed for
the classical and deep combined models and confirming, with a sample-resolved
attribution method, that the clinical variables contribute largely independent
risk information rather than refining the imaging signal. These SHAP
attributions were computed on a modest per-fold subsample and so are
qualitative rather than a basis for fine-grained ranking, but their agreement
with the classical and deep analyses reinforces that a single, biologically
coherent radiomic-plus-clinical signature drives prediction across every model
family, including the foundation model. Finally, a selected
set of influential features was back-projected to the image as voxel-level maps
in csPCa-positive cases with human-expert lesion delineations
(Supplementary Material, Figure~S2).
The displayed set retained the ADC square-root first-order total-energy map and
the recurrent DWI Imc texture features, and added a DWI LoG first-order skewness
map and a DWI square GLRLM gray-level non-uniformity map to capture complementary
diffusion-weighted first-order and run-length texture patterns.
Their salient regions co-localized with the delineated csPCa lesion rather than
being distributed uniformly through the gland; the ADC total-energy map and the
diffusion-weighted texture maps repeatedly highlighted the same focus. These
feature-resolved maps provide case-level spatial explanations directly on the
anatomy, complementing the global importance summaries.

\begin{figure}[!htbp]
\centering
\safeinclude[width=\textwidth]{fig_beeswarm.png}
\caption{\textbf{Direction and magnitude of feature effects for the best model in
each condition.} Beeswarm summaries for (A) the clinical-only Transformer--CapsNet
and (D) the dual-branch Transformer (integrated-gradient attributions), and
(B) the radiomics-only and (C) the concatenation Random Forest (SHAP values), for
the eight most important features of each model (pooled out-of-fold studies). Each
point is a study, positioned by its signed attribution and coloured by the
relative feature value (low to high); attributions to the right of zero increase
the predicted csPCa probability. In the combined models (C, D) the clinical
variables---PSA density in particular---rank among the most influential features.}
\label{fig:beeswarm}
\end{figure}

\begin{figure}[!htbp]
\centering
\safeinclude[width=\textwidth]{fig_feature_stability.png}
\caption{\textbf{Stability of in-fold feature selection.} (A) Number of features
selected in exactly $k$ of the five folds; 41 features (orange) were selected in
all folds. (B) Per-fold composition of the selected set by source sequence.
Selection is reproducible across folds and dominated by diffusion-weighted
imaging.}
\label{fig:stability}
\end{figure}

\begin{figure}[!htbp]
\centering
\safeinclude[width=\textwidth]{fig_importance_concordance.png}
\caption{\textbf{Cross-model feature-importance concordance.} (A) Normalized
permutation importance of the most informative features (union of each model's
top ten) across all six models; row labels are colored by source sequence.
(B) Pairwise overlap (top-15 Jaccard) of the most important features between
models. One diffusion-weighted GLCM feature appears in every model's top fifteen,
while overall feature-level agreement is moderate.}
\label{fig:concordance}
\end{figure}

%% ---------------------------------------------------------------
%% ---------------------------------------------------------------
%% ---------------------------------------------------------------
\section{Discussion}
\label{sec:discussion}

This study benchmarked classical machine learning against deep tabular learning
and a pretrained tabular foundation model (TabFM) on an identical
biparametric-MRI radiomic feature space for csPCa detection and quantified the
value added by routine clinical variables. Three findings stand out. First, a
classical Random Forest matched or exceeded all evaluated deep tabular
pipelines on radiomics-only data, although its specified comparison with the
Transformer was not significant after multiplicity adjustment. The numerical
ranking is consistent with the broader evidence that tree ensembles remain
strong baselines for moderate-sized tabular problems with engineered features.
Notably, the deep models were selected precisely for their stronger inductive
biases (explicit feature-interaction modeling via attention, and part--whole
representation via capsule routing)---yet these did not translate into higher
discrimination under the implemented training protocol. Second, TabFM
numerically outranked every classical and deep model in all four conditions,
and its gains over the previous best concatenation and dual-branch models were
each statistically significant after Holm adjustment, even though its
radiomics-only advantage over the Random Forest was not. Because a
random-weight control that preserved TabFM's architecture and
in-context-learning procedure while discarding its pretrained parameters
collapsed to near-chance discrimination, this advantage should be attributed to
the model's pretraining across heterogeneous tabular tasks rather than to its
architecture per se; unlike the deep tabular models evaluated here, TabFM was
not trained from scratch on this cohort, so its comparison with the other
families is architecturally informative but not a like-for-like optimization
comparison. Third, adding age, PSA,
PSA density, and prostate volume
to radiomics by simple concatenation produced a small improvement in AUROC and
AUPRC that remained significant after Holm adjustment for every model family
tested, including the classical Random Forest and, more strongly, TabFM,
whereas dual-branch fusion gave no demonstrated benefit over simple feature
concatenation among the architectures in which both were compared, classical,
deep, and TabFM alike---for TabFM, the dual-stream late fusion and simple
concatenation results were statistically indistinguishable from each other,
mirroring the pattern seen for the classical and deep architectures even though
TabFM's dual-stream mechanism was architecturally distinct from their
shared-dense-classifier fusion. The interpretability analysis added a further
dimension to these benchmark results: the
selected radiomic signature was compact and highly reproducible across folds, was
dominated by diffusion-weighted texture with a smaller but consistent ADC
contribution, and converged (most clearly on a wavelet-filtered diffusion-weighted
GLCM correlation feature that ranked among the fifteen most important variables of
every model, together with complementary first-order and texture descriptors)
across otherwise dissimilar classical and deep architectures. In the combined models the picture changed
informatively: the clinical variables, and PSA density in particular, rose to
the very top of the importance ranking, indicating that their benefit came from
largely independent risk information rather than from refining the imaging signal.
Mapping the selected high-importance features back onto the anatomy showed that
their salient regions coincided with the delineated lesion (most strikingly for
the ADC square-root total-energy map and recurrent DWI texture maps) rather than
being spread uniformly through the gland, reinforcing the biological plausibility of
the models and illustrating the feature-resolved spatial explanations this
radiomics workflow can provide.

The identity and direction of the most influential features are biologically
coherent with the pathophysiology of csPCa. Clinically significant tumours are
densely cellular, which restricts water diffusion and produces the low ADC and
elevated high b-value signal that underpin the diffusion sequences; accordingly,
the models relied heavily on diffusion-weighted descriptors, and the quantitative
ADC features surfaced by the deep models (including ADC first-order total
energy) are consistent with reduced diffusion within hypercellular tumour. The
texture descriptors that recurred across models, including gray-level
co-occurrence correlation and informational measures of
correlation, quantify the spatial orderliness and predictability of voxel
intensities: malignant tissue disrupts the regular glandular architecture of the
prostate, lowering this organized correlation structure and raising
heterogeneity, which is also reflected in the gray-level non-uniformity and
NGTDM strength descriptors and in the skewness and kurtosis of the high b-value
signal, sensitive to focal restricted-diffusion foci against benign background.
The single T2w gland-shape feature (surface-to-volume ratio) likely captures
coarse morphological correlates of disease burden. The overarching tendency was
therefore that greater textural heterogeneity and more restricted diffusion
shifted predictions toward csPCa. The clinical variables that dominated the
combined models follow established urological risk relationships: higher PSA and
PSA density, older age, and smaller prostate volume are all recognized markers of
csPCa risk, and PSA density---PSA normalized to gland volume---is a particularly
strong discriminator because it jointly encodes the higher PSA and the smaller
glands typical of significant cancer, which explains why it outranked the
individual radiomic descriptors once available to the model. That these
data-driven directions reproduce known clinical relationships strengthens
confidence that the combined model captured genuine, complementary risk signal
rather than dataset-specific artefact. These attributions should nonetheless be read
as associations rather than mechanisms, and because individual descriptors are
computed after intensity-transforming filters (square, Laplacian-of-Gaussian,
wavelet) the sign of any single filtered variant depends partly on the
transform; the robust, model-agnostic signal is the recurrent diffusion
correlation/Imc texture rather than any one coefficient, and whole-gland
extraction means the features summarize the entire gland rather than a delineated
lesion.

A calibrated csPCa risk model warrants evaluation for biopsy decision-making,
triage, and second-reader workflows. At the outer-training-derived operating
point the models
achieved a balanced sensitivity and specificity of roughly 0.75, with a high
negative predictive value (about 0.89) that is well suited to ruling out csPCa
and avoiding unnecessary biopsies; because the probabilities are calibrated, the
threshold can be shifted toward higher sensitivity where avoiding missed csPCa is
the priority, as the decision-curve analysis supports across the plausible range.
The moderate positive predictive value (about 0.53) reflects the disease
prevalence and indicates that a positive prediction warrants confirmatory work-up
rather than standalone decision-making. Because the incremental value of clinical
variables is small but real and the variables are already collected in routine
practice, a combined model is a low-cost route to improved discrimination.

The radiomics-only performance obtained here is competitive with recently published
literature, where well-tuned classical models frequently establish a high bar for csPCa
detection~\citep{cuocolo2021machine,Bao2024,Lucia2024}. However, direct performance comparisons
are often confounded by evaluation design. Many earlier radiomics studies relied on smaller
single-center cohorts or lesion-level cross-validation without patient grouping, introducing
optimistic bias and severely overestimating performance~\citep{Quality2022}. By adhering to
TRIPOD+AI guidelines with strictly leakage-safe, patient-grouped cross-validation on the
large PI-CAI cohort~\citep{collins2024tripodai,saha2023artificial}, the present study provides a conservative
internal-validation benchmark. Under the implemented protocol, the deep tabular
pipelines trained from scratch did not surpass the classical ensemble, a
pattern also reported in some external evaluations.
For example, Castillo et al.\ demonstrated that while fully automated DL achieved higher internal
cross-validation performance, classical radiomics generalized better across independent external
cohorts~\citep{Castillo2022Cancers}. The engineered inputs may already encode
much of the relevant diffusion and texture information, but the present
pipeline-level comparison cannot determine whether the remaining difference
arises from model architecture or other differences in optimization and
regularization. To bound the latter possibility, we re-ran the deep models under
an optimistically biased protocol that selected the training epoch on the
held-out fold itself (Supplementary Material, Table~S1); even with this advantage the
deep models failed to surpass the classical Random Forest in the radiomics-only
setting and gained only a small, leakage-inflated margin in the combined
conditions, indicating that the classical baseline is not an artefact of an
unfavorable deep-model training schedule. TabFM's small but significant gain
over this same classical baseline, obtained without any fold-specific gradient
training and lost entirely when its pretrained weights were replaced by random
initialization, points to a different explanation than architecture or
optimization: transferable structure learned from other tabular tasks, rather
than a training-schedule artefact, may be what a foundation model can add on
top of a strong engineered feature space. Whether this transferable prior
generalizes to other radiomic cohorts, or is specific to the statistical
structure of PI-CAI-derived features, cannot be determined from a single
internal-validation cohort.

Likewise, the predictive value added by clinical variables mirrors prior multicenter findings.
Consistent with existing risk models, PSA density emerged as the strongest routine clinical
complement to MRI~\citep{RadiomicsMultisite2023,RadiomicsNomogram2022}. While studies such as
Antol{\'i}n et al.\ reported larger discrimination gains from multimodal integration, those hybrid
models incorporated PI-RADS scores and extended clinical predictors~\citep{Antolin2026Insights}.
By deliberately omitting PI-RADS to isolate the whole-gland radiomic signal, our results
demonstrate that a parsimonious set of routine variables still adds incremental
predictive information, and that this holds even for the strongest model in the
benchmark, where TabFM's concatenation and dual-stream late-fusion results were
its best overall and statistically tied. As already noted, no elaborate fusion
strategy beat simple concatenation in any model family, so the incremental value
of the clinical variables is best captured by the parsimonious combined model.

The study has limitations. It was a retrospective analysis without an independent
external cohort; all estimates derive from the single PI-CAI collection, so
generalization to other scanners, protocols, and populations remains unverified.
Because the patient-grouped folds were not center-disjoint, acquisition-specific
signatures could be shared between training and held-out folds, and the aggregate
five-fold numbers are correspondingly optimistic; the leave-one-center-out analysis
introduced here (Section~\ref{sec:res-transport}) provides the strongest
transportability signal achievable within PI-CAI and shows that discrimination
falls by roughly $0.10$--$0.14$ AUROC when a center is genuinely held out, with every
model family---including TabFM---degrading to moderate discrimination
($\approx$0.68--0.75). This is an internal surrogate for, not a replacement of, true
external validation on an independent cohort, which remains necessary. Consistent
with this, the center-stratified analysis shows that discrimination was
substantially and significantly lower at one of the three centers (PCNN) and that
the incremental value of clinical variables was not uniform across sites, so
aggregate performance can mask clinically relevant site-level heterogeneity.
PI-RADS scores could not be recovered from the public metadata, so
the radiomic and combined models could not be benchmarked against, or combined
with, the standard radiological assessment; the present comparison is therefore
best read as a controlled comparison of model families on a fixed radiomic feature
space rather than as a positioning against the clinical standard of care.
Features were extracted from whole-gland rather than lesion-level masks, and the
masks were automatically generated, introducing potential segmentation
variability; whole-gland extraction summarizes the entire gland and may dilute
focal lesion texture, so whether the model-family ranking would hold at
lesion level is untested. PSA density was originally missing for nearly one third
of studies; it was deterministically derived from PSA and prostate volume where
possible (recovering 399 of 451 studies) and only otherwise median-imputed, and a
source-complete sensitivity analysis ($n=1049$) reproduced---and slightly
exceeded---the primary results (Section~\ref{sec:res-transport},
Table~\ref{tab:psad}), indicating that the imputation did not drive the findings.
The cohort is class-imbalanced,
and the absolute gains from clinical variables, while statistically significant,
are modest. In addition, the preliminary screening of twelve classical
algorithms and the selection of condition winners reused the evaluation cohort.
The pairwise comparisons are therefore exploratory, and their bootstrap
intervals do not capture model-selection uncertainty. Nested algorithm selection
or independent validation is needed for confirmatory inference. TabFM introduces
further, model-specific caveats: it was added to the benchmark after the
classical and deep architectures had already been selected and evaluated, so
its contrasts form a separate statistical family and it did not go through the
same twelve-algorithm screening process; its combined-modality result in the
dual condition used a logistic late-fusion of two independently in-context
TabFM streams rather than the shared-dense-classifier dual-branch architecture
used for the classical and deep models, so the ``dual'' label does not denote
the same architecture across model families; and its SHAP attributions, while
methodologically comparable to those of the other models, were computed on a
modest per-fold held-out subsample and are therefore best read as a qualitative
cross-check rather than a fine-grained ranking. The provenance and scope of TabFM's
pretraining corpus--and therefore the risk that it overlaps, even indirectly,
with public tabular benchmarks--cannot be fully audited from the outside, which
should be weighed when interpreting its advantage over models trained only on
this cohort. Future work
should include external and prospective validation,
lesion-level and zonal feature extraction, richer clinical variables,
a shared-dense-classifier dual-branch implementation of tabular foundation
models directly comparable to the classical and deep architectures, and
prospective evaluation within a clinical decision pathway~\citep{Quality2022,Lucia2024}.

%% ---------------------------------------------------------------
\section{Conclusion}
\label{sec:conclusion}

In this large, multicenter PI-CAI cohort, whole-gland radiomic features derived
from biparametric MRI provided meaningful discrimination of clinically
significant prostate cancer under multicenter internal validation. Under the
implemented training protocol, a calibrated Random Forest matched or exceeded
every deep tabular pipeline trained from scratch on radiomics-only data, and
the specified Random Forest--Transformer difference was not significant after
multiplicity adjustment; this result does not establish an architectural
advantage for classical models. A pretrained tabular foundation model, TabFM,
numerically outranked the Random Forest and every deep pipeline in all four
conditions tested, and its gains in the concatenation and combined-modality
conditions were each statistically significant after multiplicity adjustment;
a random-weight control confirmed that this advantage came from TabFM's
pretrained parameters rather than from its architecture alone. Classical tree
ensembles therefore remain essential, highly competitive reference models for
radiomics studies, and may offer a more parsimonious choice when sample size is
moderate and the feature space is structured and tabular, while pretrained
tabular foundation models emerge from this benchmark as a promising, if more
complex, alternative worth further scrutiny.

Routine clinical variables supplied complementary information across every
model family tested. Adding age, PSA, PSA density, and prostate volume to the
radiomic signature produced a modest but statistically supported improvement,
particularly in AUPRC, and this improvement was most pronounced when the
combined features were passed to TabFM, which achieved the best overall
discrimination in the study whether by simple concatenation or by a
dual-stream late fusion of separately trained radiomics-only and clinical-only
TabFM streams---the two were statistically indistinguishable from each other.
Consistent with this, a shared-dense-classifier dual-branch fusion
architecture, evaluated for the classical and deep models, likewise did not
outperform straightforward feature concatenation, so across every model family
in which more than one combination strategy was tested, added fusion
complexity conferred no demonstrated benefit over simple concatenation.
Together with the interpretability analyses, these results
indicate that diffusion-weighted texture and ADC-derived descriptors capture a
reproducible imaging phenotype, whereas PSA density and the other clinical
variables contribute additional risk information that is not fully represented
by imaging. The resulting calibrated radiomics--clinical models, whether built
on a classical ensemble or on TabFM, may be useful as a quantitative adjunct for
biopsy triage or MRI second reading, especially when high negative predictive
value is prioritized; however, they should support, rather than replace,
radiological and clinical assessment.

Clinical translation will require evaluation beyond internal cross-validation.
Independent external and prospective validation should establish robustness
across institutions, scanners, acquisition protocols, segmentation methods, and
patient populations, and should assess performance at clinically prespecified
decision thresholds. Future studies should also determine whether lesion-level
or zonal radiomics, additional clinical predictors, or longitudinal information
provide gains that justify their added complexity, and should extend the
foundation-model comparison introduced here to a shared-dense-classifier
dual-branch architecture directly comparable to the classical and deep models,
and to external cohorts, where the durability of TabFM's pretraining advantage
remains to be established. Until
such validation is available, the present work should be interpreted as a
rigorous comparative benchmark and as evidence that both a relatively simple,
calibrated, interpretable classical model and a pretrained tabular foundation
model can provide competitive, and in the latter case modestly superior, csPCa
risk estimation from routinely available biparametric MRI and clinical data.

%% ---------------------------------------------------------------
\section*{Data Availability}
The anonymized PI-CAI Public Training and Development Dataset is available
through the PI-CAI challenge under the applicable data-use terms and a
CC~BY-NC~4.0 license~\citep{saha2023artificial}. The imaging data are archived
at Zenodo (\url{https://doi.org/10.5281/zenodo.6624726}), and the maintained
annotations are available from
\url{https://github.com/DIAGNijmegen/picai_labels}. The clinical manifest and
modality-specific extracted radiomic feature tables used to construct the
analysis dataset are included in the accompanying code repository. Intermediate
combined tables, trained model files, and generated reports can be reproduced
from these resources using the provided configurations and workflow scripts.

\section*{Code Availability}
The complete analysis workflow, including radiomic feature processing, model
training, statistical evaluation, interpretability analysis, and figure
generation, is publicly available at
\url{https://github.com/jesusalzate/Radiomics-Prostate-Cancer}.

\section*{Ethics Approval}
This study was a secondary analysis of fully anonymized data released for the
PI-CAI challenge and involved no new participant recruitment, intervention, or
access to directly identifiable information. For the source PI-CAI cohort, the
institutional review boards of Radboud University Medical Center, Ziekenhuis
Groep Twente, University Medical Center Groningen, and the Norwegian University
of Science and Technology waived the requirement for informed consent for the
retrospective scientific use of anonymized clinical
data~\citep{saha2023artificial}.

\section*{Author Contributions}
\textbf{Conceptualization:} J.A.A.-G., R.T.-S., J.M.G.-G., M.I.-V.
\textbf{Methodology:} J.A.A.-G., M.A.B.-O., F.G.-G.
\textbf{Software:} J.A.A.-G., M.A.B.-O., L.M.C.P., S.L.M.
\textbf{Validation:} J.A.A.-G., M.P.-T., A.N.B., C.R.T.
\textbf{Formal analysis:} J.A.A.-G., F.G.-G.
\textbf{Investigation:} J.A.A.-G., M.A.B.-O., L.M.C.P., S.L.M.
\textbf{Resources (clinical and imaging data):} M.P.-T., A.N.B., C.R.T., J.M.O.G., M.I.-V.
\textbf{Data curation:} J.A.A.-G., M.P.-T., J.M.O.G.
\textbf{Writing -- original draft:} J.A.A.-G.
\textbf{Writing -- review \& editing:} all authors.
\textbf{Visualization:} J.A.A.-G., M.A.B.-O.
\textbf{Supervision:} R.T.-S., J.M.G.-G., M.I.-V.
\textbf{Funding acquisition:} J.A.A.-G., M.I.-V.
All authors read and approved the final manuscript.

\section*{Funding}
This work was partially supported by the Asociaci\'on Espa\~nola Contra el
C\'ancer (AECC) through the Ayuda Predoctoral AECC Valencia 2025, grant ID
PRDVA258654ALZA. The funder had no role in study design, data collection, data
analysis, data interpretation, manuscript preparation, or the decision to submit
the article for publication.

\section*{Conflicts of Interest}
The authors declare no conflicts of interest.

%% ---------------------------------------------------------------
\bibliographystyle{unsrtnat}
\bibliography{references}

\end{document}
